% !TeX root = ../main.tex
\clearpage
\section{Soundness}
\label{sec:soundness}

\subsection{Expectation and Variance}

\begin{theorem}[Expectation preservation\leanref{expectation-preservation}]
\label{thm:expectation}
Suppose $\vdash e:\tyR\E$ and $e$ is domain-safe. If
$\Expect_{\mathrm{ret}}[e]$ is defined, then the transformed program also has
positive probability of returning and a defined expectation satisfying
\[
 \Expect_{\mathrm{ret}}[\trans e]=\Expect_{\mathrm{ret}}[e].
\]
\end{theorem}
where
\begin{itemize}
    \item we write
$
 \Expect_{\mathrm{ret}}[e]\mathrel{:=}
 \frac{1}{\law e(\Real)}\int_{\Real} r\,\law e(dr)
$
for the expected output conditioned on the program successfully returning
(we assume $\law e(\Real)>0$).
    \item a program is \emph{domain-safe}\leanref{domain-safety} if,
almost surely, every executed distribution has parameters in its domain and
every executed division has a nonzero divisor.
\end{itemize}

Determinization also never increases the variance:

\begin{theorem}[Variance non-increase\leanref{variance-non-increase}]
\label{thm:variance}
Suppose $\vdash e:\tyR\E$ and $e$ is domain-safe. If
$\law e(\Real)>0$ and $\int_{\Real}r^2\,\law e(dr)<\infty$, then $\trans e$
also has positive probability of returning and a finite second moment, and
\[
 \Var_{\mathrm{ret}}[\trans e]\le\Var_{\mathrm{ret}}[e].
\]
\end{theorem}
Here $\Var_{\mathrm{ret}}$ is the variance conditioned on returning:
\[
  \Var_{\mathrm{ret}}[e]=\frac1{\law e(\Real)}\int_{\Real}
    \bigl(r-\Expect_{\mathrm{ret}}[e]\bigr)^2\,\law e(dr).
\]

The expectation theorem allows $+\infty$ and $-\infty$, but excludes undefined
expectations. \Cref{sec:traces} compares programs after fixing the $\G$ draws,
even when their global expectations are undefined.

\paragraph{When is an expectation defined?\leanref{expectation-defined}}
For a nonnegative quantity, the integral is always defined, although it may be
infinite. To integrate a signed output, separate its positive and negative parts:
\[
 r^+=\max(r,0),\qquad r^-=\max(-r,0),\qquad
 A_e=\int_{\Real}r^+\,\law e(dr),\quad
 B_e=\int_{\Real}r^-\,\law e(dr).
\]
For $q_e=\law e(\Real)>0$, we have
$\Expect_{\mathrm{ret}}[e]=(A_e-B_e)/q_e$.
If both parts are finite, the expectation is finite.
If exactly one part is infinite, the expectation is $+\infty$ or $-\infty$.
If both are infinite, the expectation is undefined: the two infinities cannot
be cancelled. In particular, having a finite expectation is equivalent to
\[
 \int_{\Real}|r|\,\law e(dr)=A_e+B_e<\infty.
\]
The absolute-value integral is nonnegative, so the question is whether it is
\emph{finite}, not whether it is defined.

For example, consider
\begin{lstlisting}
let u = uniform[G](0, 1) in
let s = 2 * bernoulli[G](1/2) - 1 in
s / u + uniform[E](-1, 1)
\end{lstlisting}
The program returns almost surely and is domain-safe: the exceptional event
$u=0$ has probability zero. Yet $\int_0^1 du/u=\infty$, and the random sign makes
both the positive and negative parts infinite, so the program's expectation is undefined. In contrast, $1/u$
alone has the defined expectation $+\infty$.

Determinization replaces the final uniform draw by zero and returns $s/u$,
whose expectation is still undefined.

\paragraph{Domain safety and missing probability.}
The following well-typed program $e$ is not domain-safe:
\begin{lstlisting}
let u = gaussian[E](0, 1) in uniform[E](0, u)
\end{lstlisting}
When $u<0$, the uniform distribution's parameters are invalid and execution gets
stuck. The program returns with probability $1/2$; conditioned on returning,
its expected output is half the mean of a standard Gaussian conditioned to be
positive. Determinization sets $u=0$ and returns the mean of
$\kw{uniform}(0,0)$, which is zero.
Thus the domain-safety assumption is necessary, because:
\[
 \Expect_{\mathrm{ret}}[e]=\frac{1}{\sqrt{2\pi}}
 \qquad\text{but}\qquad
 \Expect_{\mathrm{ret}}[\trans e]=0.
\]

\paragraph{Rejection and divergence}
Rejection, as in a failed $\kw{observe}$, contributes no output.
We model $\kw{reject}$ as an infinite loop.

Let $d_e$ denote the probability of divergence\leanref{divergence}, including rejection. It is the
limit, as $n$ grows, of the probability of executing $n$ steps and still not
being a value.

\begin{proposition}[Probability preservation]
\label{prop:probability}
Suppose $\vdash e:\tyR\E$ and $e$ is domain-safe. Then $\trans e$ is domain-safe,
and
\[
 \law e(\Real)+d_e=1\leanref{return-or-diverge},
 \qquad
 \law{\trans e}(\Real)=\law e(\Real)\leanref{output-mass}.
\]
\end{proposition}

Thus the missing output mass accounts precisely for divergence and rejection;
neither the original program nor the transformed program loses probability mass due to getting stuck.

\clearpage
\subsection{When Expectations Are Undefined}
\label{sec:traces}

Consider a program that flips a fair coin until tails, then samples uniformly
from $[-2^n,2^n]$, where $n$ is the number of heads:
\begin{lstlisting}
let sample = rec f bound =>
  if bernoulli[G](1/2) < 1 then uniform[E](-bound, bound)
  else f (2 * bound)
in sample 1
\end{lstlisting}
Fixing the $\G$ draws to $1,1,0$ (two heads, then tails) leaves the single draw
$\draw{uniform}\E{-4,4}$. After $n$ heads followed by tails, the
remaining draw is $\draw{uniform}\E{-2^n,2^n}$. Determinization replaces it
by its mean, zero.

The replacement looks right for each fixed sequence of $\G$ draws---its
$\G$ trace---but \cref{thm:expectation} gives no guarantee for this program:
its expectation is undefined. It returns almost surely; writing $N$ for the number of heads
and $X$ for its output, we have
\[
 \Pr(N=n)=2^{-n-1},\qquad
 \Expect[|X|]=\sum_{n=0}^{\infty}2^{-n-1}\frac{2^n}{2}
 =\sum_{n=0}^{\infty}\frac14=\infty.
\]
By symmetry, the positive and negative parts both have infinite expectation.
The transformed program has expectation zero.

We therefore compare the programs conditional on their $\G$ traces. For a
domain-safe program of type $\tyR\E$, determinization preserves the law of terminating $\G$ traces. For almost every such trace, the original program has a finite
conditional mean and the transformed program returns that mean. This comparison requires no
global expectation.

\paragraph{Trace semantics.}
We record the $\G$ draws alongside each returned value, giving a joint measure
$J_e$ on $\Trace\times\Real$ (\cref{fig:traces}). Forgetting the output gives
the trace measure $\trlaw e$; forgetting the trace gives the original output
measure $\law e$. These are the two marginals of $J_e$.
Only successful returns contribute mass to these measures.

% !TeX root = ../main.tex
\begin{figure}[ht]
\centering
\small
\textbf{$\G$ traces $\tau\in\Trace$}\leanref{traces}
\[
 \tau ::= []\mid(D,r)::\tau,
 \qquad
 p_{D,r}(\tau,s)=((D,r)::\tau,s).
\]
\textbf{Refined output measures}\leanref{trace-semantics}
\[
\begin{aligned}
 \joint r0&=\delta_{([],r)},\\[4pt]
 \joint{C[\draw D\E{\bar v}]}{n+1}
  &=\int_{\Real}\joint{C[r]}n\,D(\bar v)(dr),\\[4pt]
 \joint{C[\draw D\G{\bar v}]}{n+1}
  &=\int_{\Real}(p_{D,r})_*\joint{C[r]}n\,D(\bar v)(dr).
\end{aligned}
\]
\[
 J_e=\sum_{n=0}^{\infty}\joint e n,
 \qquad
 \trlaw e(A)=J_e(A\times\Real)\leanref{trace-law},
 \qquad
 \law e(B)=J_e(\Trace\times B)\leanref{trace-erasure}.
\]
\caption{Trace refinement of \cref{fig:semantics}.
$\joint e n$ records the trace and output of executions returning after exactly
$n$ steps. A $\G$ draw of value $r$ from primitive $D$ prepends $(D,r)$;
$\E$ draws and mean computations add no entry. $(p_{D,r})_*$ applies
$p_{D,r}$ to each trace/output pair in the measure. Sampling requires $\operatorname{Domain}_D(\bar v)$.
The deterministic-step and zero cases are unchanged, with $J$ in place of $\mu$.}
\label{fig:traces}
\Description{Traces record primitive names and values of G draws. The refined
semantics returns an empty trace at a real value, averages over E samples without
recording them, and prepends each G sample to the continuation trace. Summing over
termination depths gives the joint law; its marginals give traces and outputs.}
\end{figure}


\noindent\begin{minipage}{\linewidth}
In the coin-flipping example, each terminating $\G$ trace has positive probability.
A trace can also contain continuous draws:
\begin{lstlisting}
let u = uniform[G](-1, 1) in
uniform[E](0, 2) / u
\end{lstlisting}
Both programs have undefined global expectations. Yet for each fixed $u\ne0$,
the remaining expression has mean $1/u$, exactly the value returned by the
transformed program. Every individual value of $u$ has probability zero, so conditioning on a trace
cannot simply mean dividing by its probability.
\end{minipage}

We write $\law e(\,\cdot\mid\tau)$ for the output distribution conditioned on
returning and on trace $\tau$, and $\Expect[e\mid\tau]$ for its mean.
These are defined through the joint measure $J_e$:
\[
\begin{aligned}
 J_e(A\times B)&=\int_A \law e(B\mid\tau)\,\trlaw e(d\tau),\\[3pt]
 \Expect[e\mid\tau]&=\int_{\Real}r\,\law e(dr\mid\tau).
\end{aligned}
\]

\begin{theorem}[Trace soundness\leanref{trace-soundness}]
\label{thm:trace-soundness}
Suppose $\vdash e:\tyR\E$ and $e$ is domain-safe. Then $\trans e$ is domain-safe,
\[
 \trlaw{\trans e}=\trlaw e,
\]
and, for $\trlaw e$-almost every trace $\tau$,
\[
 \int_{\Real}|r|\,\law e(dr\mid\tau)<\infty,
 \qquad
 \law{\trans e}(\,\cdot\mid\tau)
 =\delta_{\Expect[e\mid\tau]}.
\]
\end{theorem}

\paragraph{Variance decomposition.}
The transformed program retains only variation between the conditional means.
The original program's output can also vary after the $\G$ trace is fixed,
due to the $\E$-samples.

\begin{theorem}[Trace variance decomposition\leanref{trace-variance}]
\label{thm:trace-variance}
Suppose $\vdash e:\tyR\E$ and $e$ is domain-safe. If
$q_e=\law e(\Real)>0$ and $\int_{\Real}r^2\,\law e(dr)<\infty$, then
\[
 \underbrace{\Var_{\mathrm{ret}}[e]}_{\text{Original variance}}
 =\underbrace{\Var_{\mathrm{ret}}[\trans e]}_{\text{Variance between traces}}
  +\underbrace{\Expect_{\tau\sim\trlaw e/q_e}
    \!\left[\Var\bigl(\law e(\,\cdot\mid\tau)\bigr)\right]}_{\text{Average variance within a trace}}.
\]
\end{theorem}

The last term averages the original program's variance after fixing a trace, using
$\trlaw e/q_e$ to condition on returning. It is nonnegative and gives exactly
the variance removed by determinization.
